% --- Archon physics formalization source begin ---
% archon:physics
% archon:covers PhyXMiniProblems/problem_phyx_mini_0932.lean
% archon:source-report reports/phyx_mini/problem_phyx_mini_0932.source.json

\chapter{Physics problem phyx\_mini\_0932}
\label{ch:PhyXMiniProblems_problem_phyx_mini_0932}

\paragraph{Problem source.}
\#\# Physical scenario

Your camping buddy has an idea for a light to go inside your tent. He happens to have a powerful (and heavy!) horseshoe magnet that he bought at a surplus store. This magnet creates a \textbackslash{}( 0.20 \textbackslash{}, \textbackslash{}text\{T\} \textbackslash{}) field between two pole tips \textbackslash{}( 10 \textbackslash{}, \textbackslash{}text\{cm\} \textbackslash{}) apart. His idea is to build the hand-cranked generator shown in figure. He thinks you can make enough current to fully light a \textbackslash{}( 1.0 \textbackslash{}, \textbackslash{}Omega \textbackslash{}) lightbulb rated at \textbackslash{}( 4.0 \textbackslash{}, \textbackslash{}text\{W\} \textbackslash{}). That's not super bright, but it should be plenty of light for routine activities in the tent.

\#\# Question

With what frequency will you have to turn the crank for the maximum current to fully light the bulb?

\#\# Answer choices

- A: A: \textbackslash{}( 4.6 \textbackslash{}times 10\textasciicircum{}2 \textbackslash{}

- B: B: \textbackslash{}( 4.1 \textbackslash{}times 10\textasciicircum{}2 \textbackslash{}

- C: C: \textbackslash{}( 2.4 \textbackslash{}times 10\textasciicircum{}2 \textbackslash{}

- D: D: \textbackslash{}( 3.1 \textbackslash{}times 10\textasciicircum{}2 \textbackslash{}

\#\# Figure caption (auxiliary; use the image as primary evidence)

The image depicts a diagram of an electrical circuit with the following elements:

1. **Semicircular Loop**: A loop with a radius labeled as 5.0 cm. It is part of the circuit and is situated within a magnetic field.

2. **Magnetic Field**: Represented by blue dots, with a strength of 0.20 T (teslas). The field is perpendicular to the plane of the semicircular loop.

3. **Crank**: Attached to the circuit on the right side, indicating a mechanism to possibly rotate or move the loop within the magnetic field.

4. **Bulb**: A symbol for a light bulb is included at the bottom of the circuit. It has specifications of 1.0 Ω (ohm) and 4.0 W (watts).

5. **Arrows**: Two circular arrows on either side suggest rotational or cyclical motion, likely related to the operation of the crank.

The setup illustrates a basic electromagnetic generator concept, where moving the loop or conducting wire within a magnetic field could induce an electric current to light the bulb.

\#\# Dataset metadata

Electromagnetic Induction; Physical Model Grounding Reasoning, Multi-Formula Reasoning

\paragraph{Recorded answer/context.}
B: B: \textbackslash{}( 4.1 \textbackslash{}times 10\textasciicircum{}2 \textbackslash{}

\paragraph{Figure/image path.}
/root/proposal\_for\_physic/science-mango/phyx\_mini\_run/phyx\_data/test\_image/932.png

\paragraph{Formalization target.}
create a compiling Lean file with sorry bodies at `PhyXMiniProblems/problem\_phyx\_mini\_0932.lean`. The Lean declarations must preserve the physical quantities, dimensions or dimensional roles, figure labels, governing-law hypotheses, and final relation expressed by this problem.
Use Mathlib/Physlib names found through LeanExplore where available. If a physics API is missing, introduce faithful local abstractions rather than scalar placeholder aliases.

\begin{theorem}[Physics formalization target]
\label{thm:physics:phyx_mini_0932:target}
\lean{PhyXMiniProblems.ProblemPhyXMini0932.requiredCrankFrequency_matches_recordedAnswerB}
\uses{def:physics:phyx-mini-0932:phyxminiproblems-problemphyxmini0932-frequencyinhertz, def:physics:phyx-mini-0932:phyxminiproblems-problemphyxmini0932-semicirculargeneratorsetup, def:physics:phyx-mini-0932:phyxminiproblems-problemphyxmini0932-matchesproblemandprimaryfigure, def:physics:phyx-mini-0932:phyxminiproblems-problemphyxmini0932-hasphysicalgeneratorparameters, def:physics:phyx-mini-0932:phyxminiproblems-problemphyxmini0932-modelsuniformfieldacrossrotatingloop, def:physics:phyx-mini-0932:phyxminiproblems-problemphyxmini0932-satisfiessemicircularloopgeometry, def:physics:phyx-mini-0932:phyxminiproblems-problemphyxmini0932-satisfiesuniformcrankrotation, def:physics:phyx-mini-0932:phyxminiproblems-problemphyxmini0932-satisfiesrotatingloopfaradaylaw, def:physics:phyx-mini-0932:phyxminiproblems-problemphyxmini0932-satisfiesratedresistivebulbcircuit, def:physics:phyx-mini-0932:phyxminiproblems-problemphyxmini0932-recordedanswerchoice, def:physics:phyx-mini-0932:phyxminiproblems-problemphyxmini0932-matchesanswerchoice}
The assigned autoformalize agent should translate this physics problem into Lean declarations in the covered file, with theorem and lemma proof bodies written as `by sorry`.
\end{theorem}
\begin{proof}
This is an autoformalization task, not a proof task. Produce statements that can later be proved by the physics prover without weakening the physical model.
\end{proof}

% --- Archon declaration topology begin ---
\section{Declaration topology}
\begin{definition}[magnetic Flux Density Dimension]
  \label{def:physics:phyx-mini-0932:phyxminiproblems-problemphyxmini0932-magneticfluxdensitydimension}
  \lean{PhyXMiniProblems.ProblemPhyXMini0932.magneticFluxDensityDimension}
  The dimension M T⁻¹ C⁻¹ of magnetic flux density (tesla)
\end{definition}
\begin{proof}
  This is the declared model definition of magnetic Flux Density Dimension.
\end{proof}
\begin{definition}[magnetic Flux Dimension]
  \label{def:physics:phyx-mini-0932:phyxminiproblems-problemphyxmini0932-magneticfluxdimension}
  \lean{PhyXMiniProblems.ProblemPhyXMini0932.magneticFluxDimension}
  The dimension M L² T⁻¹ C⁻¹ of magnetic flux (weber)
\end{definition}
\begin{proof}
  This is the declared model definition of magnetic Flux Dimension.
\end{proof}
\begin{definition}[electromotive Force Dimension]
  \label{def:physics:phyx-mini-0932:phyxminiproblems-problemphyxmini0932-electromotiveforcedimension}
  \lean{PhyXMiniProblems.ProblemPhyXMini0932.electromotiveForceDimension}
  The dimension M L² T⁻² C⁻¹ of electromotive force (volt)
\end{definition}
\begin{proof}
  This is the declared model definition of electromotive Force Dimension.
\end{proof}
\begin{definition}[electric Current Dimension]
  \label{def:physics:phyx-mini-0932:phyxminiproblems-problemphyxmini0932-electriccurrentdimension}
  \lean{PhyXMiniProblems.ProblemPhyXMini0932.electricCurrentDimension}
  The dimension C T⁻¹ of electric current (ampere)
\end{definition}
\begin{proof}
  This is the declared model definition of electric Current Dimension.
\end{proof}
\begin{definition}[electrical Resistance Dimension]
  \label{def:physics:phyx-mini-0932:phyxminiproblems-problemphyxmini0932-electricalresistancedimension}
  \lean{PhyXMiniProblems.ProblemPhyXMini0932.electricalResistanceDimension}
  The dimension M L² T⁻¹ C⁻² of electrical resistance (ohm)
\end{definition}
\begin{proof}
  This is the declared model definition of electrical Resistance Dimension.
\end{proof}
\begin{definition}[power Dimension]
  \label{def:physics:phyx-mini-0932:phyxminiproblems-problemphyxmini0932-powerdimension}
  \lean{PhyXMiniProblems.ProblemPhyXMini0932.powerDimension}
  The dimension M L² T⁻³ of power (watt)
\end{definition}
\begin{proof}
  This is the declared model definition of power Dimension.
\end{proof}
\begin{definition}[Length Magnitude]
  \label{def:physics:phyx-mini-0932:phyxminiproblems-problemphyxmini0932-lengthmagnitude}
  \lean{PhyXMiniProblems.ProblemPhyXMini0932.LengthMagnitude}
  A nonnegative, unit-independent physical length
\end{definition}
\begin{proof}
  This is the declared model definition of Length Magnitude.
\end{proof}
\begin{definition}[Area Magnitude]
  \label{def:physics:phyx-mini-0932:phyxminiproblems-problemphyxmini0932-areamagnitude}
  \lean{PhyXMiniProblems.ProblemPhyXMini0932.AreaMagnitude}
  A nonnegative, unit-independent physical area
\end{definition}
\begin{proof}
  This is the declared model definition of Area Magnitude.
\end{proof}
\begin{definition}[Magnetic Flux Density Magnitude]
  \label{def:physics:phyx-mini-0932:phyxminiproblems-problemphyxmini0932-magneticfluxdensitymagnitude}
  \lean{PhyXMiniProblems.ProblemPhyXMini0932.MagneticFluxDensityMagnitude}
  \uses{def:physics:phyx-mini-0932:phyxminiproblems-problemphyxmini0932-magneticfluxdensitydimension}
  A nonnegative, unit-independent magnetic-flux-density magnitude
\end{definition}
\begin{proof}
  This is the declared model definition of Magnetic Flux Density Magnitude.
\end{proof}
\begin{definition}[Magnetic Flux Magnitude]
  \label{def:physics:phyx-mini-0932:phyxminiproblems-problemphyxmini0932-magneticfluxmagnitude}
  \lean{PhyXMiniProblems.ProblemPhyXMini0932.MagneticFluxMagnitude}
  \uses{def:physics:phyx-mini-0932:phyxminiproblems-problemphyxmini0932-magneticfluxdimension}
  A nonnegative, unit-independent magnetic-flux amplitude
\end{definition}
\begin{proof}
  This is the declared model definition of Magnetic Flux Magnitude.
\end{proof}
\begin{definition}[Frequency Magnitude]
  \label{def:physics:phyx-mini-0932:phyxminiproblems-problemphyxmini0932-frequencymagnitude}
  \lean{PhyXMiniProblems.ProblemPhyXMini0932.FrequencyMagnitude}
  A nonnegative ordinary or angular frequency, of inverse-time dimension
\end{definition}
\begin{proof}
  This is the declared model definition of Frequency Magnitude.
\end{proof}
\begin{definition}[Electromotive Force Magnitude]
  \label{def:physics:phyx-mini-0932:phyxminiproblems-problemphyxmini0932-electromotiveforcemagnitude}
  \lean{PhyXMiniProblems.ProblemPhyXMini0932.ElectromotiveForceMagnitude}
  \uses{def:physics:phyx-mini-0932:phyxminiproblems-problemphyxmini0932-electromotiveforcedimension}
  A nonnegative, unit-independent electromotive-force magnitude
\end{definition}
\begin{proof}
  This is the declared model definition of Electromotive Force Magnitude.
\end{proof}
\begin{definition}[Electric Current Magnitude]
  \label{def:physics:phyx-mini-0932:phyxminiproblems-problemphyxmini0932-electriccurrentmagnitude}
  \lean{PhyXMiniProblems.ProblemPhyXMini0932.ElectricCurrentMagnitude}
  \uses{def:physics:phyx-mini-0932:phyxminiproblems-problemphyxmini0932-electriccurrentdimension}
  A nonnegative, unit-independent electric-current magnitude
\end{definition}
\begin{proof}
  This is the declared model definition of Electric Current Magnitude.
\end{proof}
\begin{definition}[Resistance Magnitude]
  \label{def:physics:phyx-mini-0932:phyxminiproblems-problemphyxmini0932-resistancemagnitude}
  \lean{PhyXMiniProblems.ProblemPhyXMini0932.ResistanceMagnitude}
  \uses{def:physics:phyx-mini-0932:phyxminiproblems-problemphyxmini0932-electricalresistancedimension}
  A nonnegative, unit-independent electrical resistance
\end{definition}
\begin{proof}
  This is the declared model definition of Resistance Magnitude.
\end{proof}
\begin{definition}[Power Magnitude]
  \label{def:physics:phyx-mini-0932:phyxminiproblems-problemphyxmini0932-powermagnitude}
  \lean{PhyXMiniProblems.ProblemPhyXMini0932.PowerMagnitude}
  \uses{def:physics:phyx-mini-0932:phyxminiproblems-problemphyxmini0932-powerdimension}
  A nonnegative, unit-independent rate of energy transfer
\end{definition}
\begin{proof}
  This is the declared model definition of Power Magnitude.
\end{proof}
\begin{definition}[length In Meters]
  \label{def:physics:phyx-mini-0932:phyxminiproblems-problemphyxmini0932-lengthinmeters}
  \lean{PhyXMiniProblems.ProblemPhyXMini0932.lengthInMeters}
  \uses{def:physics:phyx-mini-0932:phyxminiproblems-problemphyxmini0932-lengthmagnitude}
  Coherent-SI readout of length, in metres
\end{definition}
\begin{proof}
  This is the declared model definition of length In Meters.
\end{proof}
\begin{definition}[area In Square Meters]
  \label{def:physics:phyx-mini-0932:phyxminiproblems-problemphyxmini0932-areainsquaremeters}
  \lean{PhyXMiniProblems.ProblemPhyXMini0932.areaInSquareMeters}
  \uses{def:physics:phyx-mini-0932:phyxminiproblems-problemphyxmini0932-areamagnitude}
  Coherent-SI readout of area, in square metres
\end{definition}
\begin{proof}
  This is the declared model definition of area In Square Meters.
\end{proof}
\begin{definition}[magnetic Flux Density In Teslas]
  \label{def:physics:phyx-mini-0932:phyxminiproblems-problemphyxmini0932-magneticfluxdensityinteslas}
  \lean{PhyXMiniProblems.ProblemPhyXMini0932.magneticFluxDensityInTeslas}
  \uses{def:physics:phyx-mini-0932:phyxminiproblems-problemphyxmini0932-magneticfluxdensitymagnitude}
  Coherent-SI readout of magnetic flux density, in teslas
\end{definition}
\begin{proof}
  This is the declared model definition of magnetic Flux Density In Teslas.
\end{proof}
\begin{definition}[magnetic Flux In Webers]
  \label{def:physics:phyx-mini-0932:phyxminiproblems-problemphyxmini0932-magneticfluxinwebers}
  \lean{PhyXMiniProblems.ProblemPhyXMini0932.magneticFluxInWebers}
  \uses{def:physics:phyx-mini-0932:phyxminiproblems-problemphyxmini0932-magneticfluxmagnitude}
  Coherent-SI readout of magnetic flux, in webers
\end{definition}
\begin{proof}
  This is the declared model definition of magnetic Flux In Webers.
\end{proof}
\begin{definition}[frequency In Hertz]
  \label{def:physics:phyx-mini-0932:phyxminiproblems-problemphyxmini0932-frequencyinhertz}
  \lean{PhyXMiniProblems.ProblemPhyXMini0932.frequencyInHertz}
  \uses{def:physics:phyx-mini-0932:phyxminiproblems-problemphyxmini0932-frequencymagnitude}
  Coherent-SI readout of ordinary frequency, in hertz
\end{definition}
\begin{proof}
  This is the declared model definition of frequency In Hertz.
\end{proof}
\begin{definition}[angular Frequency In Radians Per Second]
  \label{def:physics:phyx-mini-0932:phyxminiproblems-problemphyxmini0932-angularfrequencyinradianspersecond}
  \lean{PhyXMiniProblems.ProblemPhyXMini0932.angularFrequencyInRadiansPerSecond}
  \uses{def:physics:phyx-mini-0932:phyxminiproblems-problemphyxmini0932-frequencymagnitude}
  Coherent-SI readout of angular frequency, in radians per second
\end{definition}
\begin{proof}
  This is the declared model definition of angular Frequency In Radians Per Second.
\end{proof}
\begin{definition}[electromotive Force In Volts]
  \label{def:physics:phyx-mini-0932:phyxminiproblems-problemphyxmini0932-electromotiveforceinvolts}
  \lean{PhyXMiniProblems.ProblemPhyXMini0932.electromotiveForceInVolts}
  \uses{def:physics:phyx-mini-0932:phyxminiproblems-problemphyxmini0932-electromotiveforcemagnitude}
  Coherent-SI readout of electromotive force, in volts
\end{definition}
\begin{proof}
  This is the declared model definition of electromotive Force In Volts.
\end{proof}
\begin{definition}[electric Current In Amperes]
  \label{def:physics:phyx-mini-0932:phyxminiproblems-problemphyxmini0932-electriccurrentinamperes}
  \lean{PhyXMiniProblems.ProblemPhyXMini0932.electricCurrentInAmperes}
  \uses{def:physics:phyx-mini-0932:phyxminiproblems-problemphyxmini0932-electriccurrentmagnitude}
  Coherent-SI readout of electric current, in amperes
\end{definition}
\begin{proof}
  This is the declared model definition of electric Current In Amperes.
\end{proof}
\begin{definition}[resistance In Ohms]
  \label{def:physics:phyx-mini-0932:phyxminiproblems-problemphyxmini0932-resistanceinohms}
  \lean{PhyXMiniProblems.ProblemPhyXMini0932.resistanceInOhms}
  \uses{def:physics:phyx-mini-0932:phyxminiproblems-problemphyxmini0932-resistancemagnitude}
  Coherent-SI readout of electrical resistance, in ohms
\end{definition}
\begin{proof}
  This is the declared model definition of resistance In Ohms.
\end{proof}
\begin{definition}[power In Watts]
  \label{def:physics:phyx-mini-0932:phyxminiproblems-problemphyxmini0932-powerinwatts}
  \lean{PhyXMiniProblems.ProblemPhyXMini0932.powerInWatts}
  \uses{def:physics:phyx-mini-0932:phyxminiproblems-problemphyxmini0932-powermagnitude}
  Coherent-SI readout of power, in watts
\end{definition}
\begin{proof}
  This is the declared model definition of power In Watts.
\end{proof}
\begin{definition}[Figure Object]
  \label{def:physics:phyx-mini-0932:phyxminiproblems-problemphyxmini0932-figureobject}
  \lean{PhyXMiniProblems.ProblemPhyXMini0932.FigureObject}
  Physical components visibly present in image 932.png
\end{definition}
\begin{proof}
  This is the declared model definition of Figure Object.
\end{proof}
\begin{definition}[Figure Label]
  \label{def:physics:phyx-mini-0932:phyxminiproblems-problemphyxmini0932-figurelabel}
  \lean{PhyXMiniProblems.ProblemPhyXMini0932.FigureLabel}
  Textual or symbolic labels printed in image 932.png
\end{definition}
\begin{proof}
  This is the declared model definition of Figure Label.
\end{proof}
\begin{definition}[Normal Field Glyph]
  \label{def:physics:phyx-mini-0932:phyxminiproblems-problemphyxmini0932-normalfieldglyph}
  \lean{PhyXMiniProblems.ProblemPhyXMini0932.NormalFieldGlyph}
  The dot glyph in the figure denotes a field directed out of the page
\end{definition}
\begin{proof}
  This is the declared model definition of Normal Field Glyph.
\end{proof}
\begin{definition}[Diagram Direction]
  \label{def:physics:phyx-mini-0932:phyxminiproblems-problemphyxmini0932-diagramdirection}
  \lean{PhyXMiniProblems.ProblemPhyXMini0932.DiagramDirection}
  Named directions needed to transcribe the supplied front-view figure
\end{definition}
\begin{proof}
  This is the declared model definition of Diagram Direction.
\end{proof}
\begin{definition}[Loop Rotation Axis]
  \label{def:physics:phyx-mini-0932:phyxminiproblems-problemphyxmini0932-looprotationaxis}
  \lean{PhyXMiniProblems.ProblemPhyXMini0932.LoopRotationAxis}
  The rotation axis selected by the two axle points in the diagram
\end{definition}
\begin{proof}
  This is the declared model definition of Loop Rotation Axis.
\end{proof}
\begin{definition}[Generator Conductor Model]
  \label{def:physics:phyx-mini-0932:phyxminiproblems-problemphyxmini0932-generatorconductormodel}
  \lean{PhyXMiniProblems.ProblemPhyXMini0932.GeneratorConductorModel}
  The conducting geometry idealized by this school-physics model
\end{definition}
\begin{proof}
  This is the declared model definition of Generator Conductor Model.
\end{proof}
\begin{definition}[Semicircular Generator Figure]
  \label{def:physics:phyx-mini-0932:phyxminiproblems-problemphyxmini0932-semicirculargeneratorfigure}
  \lean{PhyXMiniProblems.ProblemPhyXMini0932.SemicircularGeneratorFigure}
  \uses{def:physics:phyx-mini-0932:phyxminiproblems-problemphyxmini0932-figureobject, def:physics:phyx-mini-0932:phyxminiproblems-problemphyxmini0932-figurelabel, def:physics:phyx-mini-0932:phyxminiproblems-problemphyxmini0932-normalfieldglyph, def:physics:phyx-mini-0932:phyxminiproblems-problemphyxmini0932-diagramdirection, def:physics:phyx-mini-0932:phyxminiproblems-problemphyxmini0932-looprotationaxis}
  Typed presentation data copied from the primary raster. It contains no crank frequency or answer-choice value
\end{definition}
\begin{proof}
  This is the declared model definition of Semicircular Generator Figure.
\end{proof}
\begin{definition}[Semicircular Generator Setup]
  \label{def:physics:phyx-mini-0932:phyxminiproblems-problemphyxmini0932-semicirculargeneratorsetup}
  \lean{PhyXMiniProblems.ProblemPhyXMini0932.SemicircularGeneratorSetup}
  \uses{def:physics:phyx-mini-0932:phyxminiproblems-problemphyxmini0932-lengthmagnitude, def:physics:phyx-mini-0932:phyxminiproblems-problemphyxmini0932-areamagnitude, def:physics:phyx-mini-0932:phyxminiproblems-problemphyxmini0932-magneticfluxdensitymagnitude, def:physics:phyx-mini-0932:phyxminiproblems-problemphyxmini0932-magneticfluxmagnitude, def:physics:phyx-mini-0932:phyxminiproblems-problemphyxmini0932-frequencymagnitude, def:physics:phyx-mini-0932:phyxminiproblems-problemphyxmini0932-electromotiveforcemagnitude, def:physics:phyx-mini-0932:phyxminiproblems-problemphyxmini0932-electriccurrentmagnitude, def:physics:phyx-mini-0932:phyxminiproblems-problemphyxmini0932-resistancemagnitude, def:physics:phyx-mini-0932:phyxminiproblems-problemphyxmini0932-powermagnitude, def:physics:phyx-mini-0932:phyxminiproblems-problemphyxmini0932-generatorconductormodel, def:physics:phyx-mini-0932:phyxminiproblems-problemphyxmini0932-semicirculargeneratorfigure}
  Independent physical quantities describing the generator and bulb. The crank frequency is an unknown observable, not a definition made from the requested answer. Peak emf, current, flux, and power are likewise independent fields related only by the law predicates below
\end{definition}
\begin{proof}
  This is the declared model definition of Semicircular Generator Setup.
\end{proof}
\begin{definition}[Matches Problem And Primary Figure]
  \label{def:physics:phyx-mini-0932:phyxminiproblems-problemphyxmini0932-matchesproblemandprimaryfigure}
  \lean{PhyXMiniProblems.ProblemPhyXMini0932.MatchesProblemAndPrimaryFigure}
  \uses{def:physics:phyx-mini-0932:phyxminiproblems-problemphyxmini0932-lengthinmeters, def:physics:phyx-mini-0932:phyxminiproblems-problemphyxmini0932-magneticfluxdensityinteslas, def:physics:phyx-mini-0932:phyxminiproblems-problemphyxmini0932-resistanceinohms, def:physics:phyx-mini-0932:phyxminiproblems-problemphyxmini0932-powerinwatts, def:physics:phyx-mini-0932:phyxminiproblems-problemphyxmini0932-semicirculargeneratorsetup}
  Numerical prose data and literal qualitative evidence from image 932.png. The unknown crank frequency and all equivalent formulas for it are absent
\end{definition}
\begin{proof}
  This is the declared model definition of Matches Problem And Primary Figure.
\end{proof}
\begin{definition}[Has Physical Generator Parameters]
  \label{def:physics:phyx-mini-0932:phyxminiproblems-problemphyxmini0932-hasphysicalgeneratorparameters}
  \lean{PhyXMiniProblems.ProblemPhyXMini0932.HasPhysicalGeneratorParameters}
  \uses{def:physics:phyx-mini-0932:phyxminiproblems-problemphyxmini0932-lengthinmeters, def:physics:phyx-mini-0932:phyxminiproblems-problemphyxmini0932-areainsquaremeters, def:physics:phyx-mini-0932:phyxminiproblems-problemphyxmini0932-magneticfluxdensityinteslas, def:physics:phyx-mini-0932:phyxminiproblems-problemphyxmini0932-magneticfluxinwebers, def:physics:phyx-mini-0932:phyxminiproblems-problemphyxmini0932-frequencyinhertz, def:physics:phyx-mini-0932:phyxminiproblems-problemphyxmini0932-angularfrequencyinradianspersecond, def:physics:phyx-mini-0932:phyxminiproblems-problemphyxmini0932-electromotiveforceinvolts, def:physics:phyx-mini-0932:phyxminiproblems-problemphyxmini0932-electriccurrentinamperes, def:physics:phyx-mini-0932:phyxminiproblems-problemphyxmini0932-resistanceinohms, def:physics:phyx-mini-0932:phyxminiproblems-problemphyxmini0932-powerinwatts, def:physics:phyx-mini-0932:phyxminiproblems-problemphyxmini0932-semicirculargeneratorsetup}
  Positivity and nondegeneracy conditions for the operating generator
\end{definition}
\begin{proof}
  This is the declared model definition of Has Physical Generator Parameters.
\end{proof}
\begin{definition}[Models Uniform Field Across Rotating Loop]
  \label{def:physics:phyx-mini-0932:phyxminiproblems-problemphyxmini0932-modelsuniformfieldacrossrotatingloop}
  \lean{PhyXMiniProblems.ProblemPhyXMini0932.ModelsUniformFieldAcrossRotatingLoop}
  \uses{def:physics:phyx-mini-0932:phyxminiproblems-problemphyxmini0932-magneticfluxdensityinteslas, def:physics:phyx-mini-0932:phyxminiproblems-problemphyxmini0932-semicirculargeneratorsetup}
  The physical field is uniform in magnitude over the region between the pole tips, and the rotating semicircular conductor remains in that active region
\end{definition}
\begin{proof}
  This is the declared model definition of Models Uniform Field Across Rotating Loop.
\end{proof}
\begin{definition}[Satisfies Semicircular Loop Geometry]
  \label{def:physics:phyx-mini-0932:phyxminiproblems-problemphyxmini0932-satisfiessemicircularloopgeometry}
  \lean{PhyXMiniProblems.ProblemPhyXMini0932.SatisfiesSemicircularLoopGeometry}
  \uses{def:physics:phyx-mini-0932:phyxminiproblems-problemphyxmini0932-lengthinmeters, def:physics:phyx-mini-0932:phyxminiproblems-problemphyxmini0932-areainsquaremeters, def:physics:phyx-mini-0932:phyxminiproblems-problemphyxmini0932-semicirculargeneratorsetup}
  The area enclosed by a semicircle of radius r is π r² / 2
\end{definition}
\begin{proof}
  This is the declared model definition of Satisfies Semicircular Loop Geometry.
\end{proof}
\begin{definition}[Satisfies Uniform Crank Rotation]
  \label{def:physics:phyx-mini-0932:phyxminiproblems-problemphyxmini0932-satisfiesuniformcrankrotation}
  \lean{PhyXMiniProblems.ProblemPhyXMini0932.SatisfiesUniformCrankRotation}
  \uses{def:physics:phyx-mini-0932:phyxminiproblems-problemphyxmini0932-frequencyinhertz, def:physics:phyx-mini-0932:phyxminiproblems-problemphyxmini0932-angularfrequencyinradianspersecond, def:physics:phyx-mini-0932:phyxminiproblems-problemphyxmini0932-semicirculargeneratorsetup}
  Uniform crank rotation relates angular and ordinary frequency by ω=2πf
\end{definition}
\begin{proof}
  This is the declared model definition of Satisfies Uniform Crank Rotation.
\end{proof}
\begin{definition}[Satisfies Rotating Loop Faraday Law]
  \label{def:physics:phyx-mini-0932:phyxminiproblems-problemphyxmini0932-satisfiesrotatingloopfaradaylaw}
  \lean{PhyXMiniProblems.ProblemPhyXMini0932.SatisfiesRotatingLoopFaradayLaw}
  \uses{def:physics:phyx-mini-0932:phyxminiproblems-problemphyxmini0932-areainsquaremeters, def:physics:phyx-mini-0932:phyxminiproblems-problemphyxmini0932-magneticfluxdensityinteslas, def:physics:phyx-mini-0932:phyxminiproblems-problemphyxmini0932-magneticfluxinwebers, def:physics:phyx-mini-0932:phyxminiproblems-problemphyxmini0932-angularfrequencyinradianspersecond, def:physics:phyx-mini-0932:phyxminiproblems-problemphyxmini0932-electromotiveforceinvolts, def:physics:phyx-mini-0932:phyxminiproblems-problemphyxmini0932-semicirculargeneratorsetup}
  For the single rotating semicircular loop, the flux amplitude is B A, and the peak form of Faraday's law is E\_max = Φ\_max ω. These are generic generator laws and do not state a numerical crank frequency
\end{definition}
\begin{proof}
  This is the declared model definition of Satisfies Rotating Loop Faraday Law.
\end{proof}
\begin{definition}[Satisfies Rated Resistive Bulb Circuit]
  \label{def:physics:phyx-mini-0932:phyxminiproblems-problemphyxmini0932-satisfiesratedresistivebulbcircuit}
  \lean{PhyXMiniProblems.ProblemPhyXMini0932.SatisfiesRatedResistiveBulbCircuit}
  \uses{def:physics:phyx-mini-0932:phyxminiproblems-problemphyxmini0932-electromotiveforceinvolts, def:physics:phyx-mini-0932:phyxminiproblems-problemphyxmini0932-electriccurrentinamperes, def:physics:phyx-mini-0932:phyxminiproblems-problemphyxmini0932-resistanceinohms, def:physics:phyx-mini-0932:phyxminiproblems-problemphyxmini0932-powerinwatts, def:physics:phyx-mini-0932:phyxminiproblems-problemphyxmini0932-semicirculargeneratorsetup}
  The bulb is the only appreciable resistance. At the requested operating point, Ohm's law and Joule heating make the peak bulb power equal its rating
\end{definition}
\begin{proof}
  This is the declared model definition of Satisfies Rated Resistive Bulb Circuit.
\end{proof}
\begin{lemma}[required Crank Frequency exact]
  \label{lem:physics:phyx-mini-0932:phyxminiproblems-problemphyxmini0932-requiredcrankfrequency-exact}
  \lean{PhyXMiniProblems.ProblemPhyXMini0932.requiredCrankFrequency_exact}
  \uses{def:physics:phyx-mini-0932:phyxminiproblems-problemphyxmini0932-frequencyinhertz, def:physics:phyx-mini-0932:phyxminiproblems-problemphyxmini0932-semicirculargeneratorsetup, def:physics:phyx-mini-0932:phyxminiproblems-problemphyxmini0932-matchesproblemandprimaryfigure, def:physics:phyx-mini-0932:phyxminiproblems-problemphyxmini0932-hasphysicalgeneratorparameters, def:physics:phyx-mini-0932:phyxminiproblems-problemphyxmini0932-modelsuniformfieldacrossrotatingloop, def:physics:phyx-mini-0932:phyxminiproblems-problemphyxmini0932-satisfiessemicircularloopgeometry, def:physics:phyx-mini-0932:phyxminiproblems-problemphyxmini0932-satisfiesuniformcrankrotation, def:physics:phyx-mini-0932:phyxminiproblems-problemphyxmini0932-satisfiesrotatingloopfaradaylaw, def:physics:phyx-mini-0932:phyxminiproblems-problemphyxmini0932-satisfiesratedresistivebulbcircuit}
  The exact school-physics result before answer-choice rounding is 4000 / π² Hz
\end{lemma}
\begin{proof}
  The conclusion follows from the governing relations and figure readouts named in the statement of required Crank Frequency exact.
\end{proof}
\begin{definition}[Answer Choice]
  \label{def:physics:phyx-mini-0932:phyxminiproblems-problemphyxmini0932-answerchoice}
  \lean{PhyXMiniProblems.ProblemPhyXMini0932.AnswerChoice}
  Labels of the four answer choices supplied with the problem
\end{definition}
\begin{proof}
  This is the declared model definition of Answer Choice.
\end{proof}
\begin{definition}[frequency In Hertz]
  \label{def:physics:phyx-mini-0932:phyxminiproblems-problemphyxmini0932-answerchoice-frequencyinhertz}
  \lean{PhyXMiniProblems.ProblemPhyXMini0932.AnswerChoice.frequencyInHertz}
  \uses{def:physics:phyx-mini-0932:phyxminiproblems-problemphyxmini0932-frequencyinhertz, def:physics:phyx-mini-0932:phyxminiproblems-problemphyxmini0932-answerchoice}
  Frequency in hertz printed beside each answer label
\end{definition}
\begin{proof}
  This is the declared model definition of frequency In Hertz.
\end{proof}
\begin{definition}[recorded Answer Choice]
  \label{def:physics:phyx-mini-0932:phyxminiproblems-problemphyxmini0932-recordedanswerchoice}
  \lean{PhyXMiniProblems.ProblemPhyXMini0932.recordedAnswerChoice}
  \uses{def:physics:phyx-mini-0932:phyxminiproblems-problemphyxmini0932-answerchoice}
  The answer label recorded in the supplied dataset metadata
\end{definition}
\begin{proof}
  This is the declared model definition of recorded Answer Choice.
\end{proof}
\begin{definition}[Matches Answer Choice]
  \label{def:physics:phyx-mini-0932:phyxminiproblems-problemphyxmini0932-matchesanswerchoice}
  \lean{PhyXMiniProblems.ProblemPhyXMini0932.MatchesAnswerChoice}
  \uses{def:physics:phyx-mini-0932:phyxminiproblems-problemphyxmini0932-frequencymagnitude, def:physics:phyx-mini-0932:phyxminiproblems-problemphyxmini0932-frequencyinhertz, def:physics:phyx-mini-0932:phyxminiproblems-problemphyxmini0932-answerchoice}
  Agreement with a frequency displayed to the nearest 10 Hz: the exact value lies within half of one displayed 10 Hz step
\end{definition}
\begin{proof}
  This is the declared model definition of Matches Answer Choice.
\end{proof}
% --- Archon declaration topology end ---
% --- Archon physics formalization source end ---
